\section{Related Work} 
In this section, we briefly review existing studies on LLM-based ranking and utility in  Retrieval-Augmented Generation (RAG).
% \subsection{Background} 
% Information retrieval (IR) typically comprises multistage ranking pipelines: retrieval and ranking \cite{guo2022semantic}. 
% Formally, given a query $q$ and a large corpus $\mathcal{C}$ containing numerous documents, the objective of the retrieval stage is to retrieve a list of potentially relevant documents $D = \{d_1, d_2, ..., d_N\}$ from $\mathcal{C}$. 
% Dense retrieval (DR) models are currently the mainstream retrieval method, which encodes the query and documents into dense embeddings, following the dual encoder paradigm. 
% The underpinning backbone of DR has evolved, progressing from encoder-based models \cite{xiao2022retromae, karpukhin2020dense, gao2021condenser, ma2021prop, zhang2025leveraging} and encoder-decoder models \cite{ni2021large} to the contemporary paradigm dominated by decoder-only large language models (LLMs) \cite{zhang2025unleashing, ma2024fine, behnamghader2024llm2vec, li2024llama2vec, springer2024repetition, zeng2025scaling}. 
% Critically, the integration of these large-scale models has further yielded substantial enhancements in retrieval performance. 
% Subsequently, the goal of the ranking stage is to further refine the ordering of the documents within the retrieved list $D$ to optimize ranking performance. 
% Next, we briefly review some existing studies on ranking and Retrieval-Augmented Generation (RAG).  


\subsection{LLM-based Ranking} 
Ranking models frequently utilize a cross-encoder architecture, enabling fine-grained, token-level interactions between the query and the document text \cite{guo2019matchzoo}. 
Ranking approaches are categorized based on how they model the relationship between a query and documents: pointwise \cite{nogueira2019multi, nogueira2020document}, pairwise \cite{pradeep2021expando, burges2005learning}, and listwise \cite{cao2007learning, wang2018lambdaloss, gao2021rethink, zhuang2023rankt5}.  
In the era of pre-trained language models (PLMs), owing to input length constraints, listwise ranking methods are typically implemented using pointwise inputs during training while employing listwise loss functions. 
This approach cannot be considered authentic listwise learning, as it inherently fails to capture query-document interactions during inference. 
The widespread adoption of Large Language Models (LLMs) has spurred significant advancements within the ranking domain, especially zero-shot listwise ranking \cite{liu2025leveraging, reddy2024first, pradeep2023rankzephyr, ren2025self}. 
RankVicuna \cite{pradeep2023rankvicuna} and RankGPT \cite{sun2023chatgpt} pioneered the application of using LLMs for zero-shot listwise ranking. 
RankZephyr \cite{pradeep2023rankzephyr} explored distilling the powerful listwise ranking capabilities of high-performing LLMs into smaller, more efficient models. 
Currently, the primary distillation signal involves relevance ranking sequence generated by LLM. 
Unlike these works, this work explores novel methods to distill the utility-based selection capability of LLMs into smaller models. 


\subsection{Utility in RAG}
Retrieval-Augmented Generation (RAG) typically involves two stages: retrieval and generation. While relevance serves as the primary optimization objective for the retriever, the ultimate goal in RAG shifts towards optimizing end-task question answering (QA) performance using outputs from effective retrievers, with less direct emphasis on retrieval metrics themselves \cite{zhang2025leveraging}. Consequently, significant research has focused on the importance of retrieval utility – the measure of a passage's contribution to generating a correct answer – within RAG systems \cite{zhang2024iterative, zhang2025leveraging, zhang2024large, shi2023replug, izacard2023atlas}. \citet{zhang2024large} was the first to propose directly using large language models (LLMs) to judge the utility of retrieved passages, demonstrating that LLM-generated pseudo-answers can effectively assist utility assessment. \citet{zhang2024iterative} further introduced an approach based on Schutz's theory of relevance to iteratively enhance LLM-based utility judgments. However, these methods are constrained by their reliance on utility judgments for only a small number of retrieval candidates per query. 
In contrast, \citet{shi2023replug} proposed quantifying utility score probabilistically, using either the likelihood of the ground-truth answer or performance differences on downstream tasks. This probabilistic approach, however, relies heavily on the availability of ground-truth answers. 
In this work, we extend utility-based selection to a significantly larger scale of retrieval candidates.


